\documentclass[12pt]{article}

\usepackage[utf8]{inputenc}
\usepackage[T1]{fontenc}
\usepackage{geometry}
\usepackage{hyperref}

\geometry{margin=2.5cm}

\title{Terrae Novae 2030+ Strategy Roadmap}
\author{ESA}
\date{2022}

\begin{document}

\maketitle

\begin{abstract}
\begin{itemize}
\item The Human Spaceflight and Exploration Science Advisory Committee (HESAC) looks forward to the vision outlined in Terrae Novae 2030+ strategy, focused on expanding exploitation of Low Earth Orbit (LEO) and a thorough investigation of the Moon and Mars.
\item It is expected that these capabilities will secure the ESA ambition of “European oots on the Moon”, i.e. European science and technology together with a first European astronaut on the Moon surface as soon as 2030.
\item At CM22, specific decisions will be required to ensure long-term European capabilities (e.g. in LEO) and to prepare the next steps in deep space (e.g. for lunar surface exploration and preparing for humans to Mars).
\end{itemize}
\end{abstract}

\section{Fonte}

\begin{itemize}
\item \textbf{Tipo}: pdf
\item \textbf{Título}: Terrae Novae 2030+ Strategy Roadmap
\item \textbf{Autor / organização}: ESA
\item \textbf{Localização}: \url{https://esamultimedia.esa.int/docs/HRE/TerraenovaeStrategyRoadmap.pdf}
\end{itemize}

\section{Texto da Fonte}

% texto original limpo da fonte, incluído no corpo do artigo conforme pedido no enunciado
Terrae Novae 2030+ Strategy Roadmap June 2022

The mission of the Terrae Novae exploration programme is to lead Europe’s human journey into the Solar system using robots as precursors and scouts, and to return the benefits of exploration back to society. Terrae Novae has the literal meaning of the ‘New Worlds’ that encompasses the three ESA exploration destinations: Low Earth Orbit (LEO), Moon and Mars. It evokes the spirit of new discoveries, new ambitions, new science, new inspiration, and new challenges. It symbolises the constant quest for technological, process and procurement innovations that result in new and better ways to deliver the programme. It also reflects the aspiration to actively reach out to new partners from beyond the space sector and enlarge the space ecosystem to the commercial sphere. European Space Agency Foreword – by Josef Aschbacher, Director General

At the beginning of this decade, space exploration is at an unprecedented crossroads. Space exploration is unquestionably an investment for future prosperity. It demonstrably generates high quality jobs and immediate economic return. Exploration science and technologies are an accelerator of sustainable development which are already generating innovative solutions that make life on Earth more productive, cleaner, and sustainable. Space exploration is also inherently exciting, inspiring and motivating - especially for the next generation. Many times, I have witnessed for myself the extraordinary impact our ESA astronauts have upon people, both young and old. It is therefore not surprising that leading countries are accelerating investments, while emerging countries and private investors are defining and implementing their own plans. All are on a course to define their posture in exploration through strong strategies supported by visible actions. I am therefore delighted that exciting but achievable goals have been elaborated in the ESA Terrae Novae 2030+ strategy roadmap. If implemented, this roadmap will deliver benefits to all the citizens of our Member States. Undeniably, it is an exploration plan for an ambitious Europe, and one that is fully aligned with our Agenda 20251. Its main goal is to provide a lighthouse that illuminates for Europe’s decision makers a long-term vision which reaches beyond our current exploration programme and its many activities and achievements. While always delivering measurable benefits for society today and tomorrow, the top objectives are threefold: 1. to create new opportunities in Low Earth Orbit for a sustained European presence in the post-ISS era, 2. to enable the first European to explore the Moon’s surface by 2030 as a step towards sustainable lunar exploration in the 2030’s, 3. to prepare the horizon goal of Europe being part of the first human mission to Mars. How can this strategy roadmap be used? I see many purposes. When embarking on an ambitious and challenging journey, having a good roadmap is always recommended. This Terrae Novae 2030+ strategy roadmap is a flexible instrument with options to tune decision-making, considering the evolving political and programmatic landscape, and the level of ambition and affordability of our Member States. Already at the ESA Council at Ministerial level in November 2022, decisions are required to ensure long-term European presence in Low Earth Orbit, to prepare the next steps for lunar surface exploration, and to plan intermediate steps towards an eventual human Mars mission. This roadmap supports these decisions by identifying candidate new missions and related technologies that could be prepared within the Terrae Novae programme in the 2025-2030 timeframe. When viewed from a strategic perspective, the roadmap sets out a consolidated proposal for use by all European stakeholders - governments, space agencies, the science community, and industry including the non-space sector. It thus provides a narrative needed for political decision makers and also for taxpayers’ appreciation concerning what is already planned and what is possible in the future. I believe it also sends a message towards our valued international partners that Europe has both a vision and a direction of travel. Last, but not least, I hope the document can serve as a useful reference for the work of the High-level Advisory Group on Human Space Exploration for Europe, which was mandated during the European Space Summit on 16 February 2022 in Toulouse, France. More than any other space activity, space exploration offers a unique blend of curiosity and opportunity - the curiosity to venture into the unknown in search of new horizons and new knowledge; and the opportunity to return to society the many benefits of making the journey. I now invite our political decision-makers to define Europe’s level of ambition so that ESA, together with all its stakeholders, can translate this strategy roadmap into reality. Josef Aschbacher Director General European Space Agency June 2022 Exploring the Solar System – by the Human spaceflight and Exploration Science Advisory Committee (HESAC)

The Human Spaceflight and Exploration Science Advisory Committee (HESAC) looks forward to the vision outlined in Terrae Novae 2030+ strategy, focused on expanding exploitation of Low Earth Orbit (LEO) and a thorough investigation of the Moon and Mars. The strategy sets a robust basis for the continuation of the previous programme and capitalizes on current achievements. It also offers new opportunities and resilience elements that are essential given the current political climate. Indeed, this new strategic plan is based on a scientifically diverse and innovative outlook that reflects our growing European aspirations to provide Europe with the necessary autonomy and, in areas of our established expertise, potential leadership positions for all three destinations in the coming years and even decades. The scope of this new strategy will expand our knowledge of the Solar System and its early history, galvanised by international collaboration, with our traditional partners such as NASA and others, and the emerging commercial sector. This will also allow for the acceleration of research developments and technological breakthroughs. All these advances are offered by the exploitation of LEO and cis-lunar infrastructures (with the critical contribution to the Orion vehicle and the cis-lunar Gateway for example), the robotic and sustained human access to the surface of the Moon (with the European Large Logistics Lander) and the long-term revisit of Mars for comprehensive state-of-the-art scientific studies and sample return (after ExoMars/Trace Gas Orbiter and then with Rosalind Franklin and Mars Sample Return, followed by human exploration in the early 2040s). The development of new robotic techniques, together with human-assisted robotic instruments, will yield precious scientific data from important planetary locations. In addition, the prospect of analysing samples returned from the Moon and Mars in sophisticated terrestrial laboratories will enhance our understanding of the physical and chemical processes underlying the origin and evolution of our immediate environment (our own planet and its neighbours). The Terrae Novae 2030+ strategy fosters strong collaborations with a wide range of international partners in the development of exploration missions and projects, allowing for ESA to be established in privileged roles promoting European science and very challenging engineering and corresponding translational biomedical benefits foreseen in the next phase of exploration. This strategic plan provides the means of achieving ESA’s ambitious exploration goals and acts as a focus for public engagement, advocacy and increased commercial opportunities in this long-term international endeavour on the basis of new capabilities and missions envisioned by ESA and its partners. The collaboration among ESA’s science, technology and space transportation directorates is an important asset for the strategy, as is the reliability that ESA offers as an international partner. Exploring the Solar System – by the Human spaceflight and Exploration Science Advisory Committee

As recent events have shown, the geopolitical context can unexpectedly become unstable. Consequently, historical international cooperation, even for the highly emblematic and peaceful activity of robotic and human space exploration, can suddenly be called into question. It is a reminder that political tensions and change are difficult to predict, and that resilience is of the essence. More isolationism and economic protectionism are unfortunately a trend, and new potential superpower confrontations and remodelling of alliances can still be expected. In addition, bottom-up disruptive technologies and approaches from the private sector are now triggered from a purely commercial angle and can profoundly reshuffle the cards at any time as new(space) enablers of exploration emerge. Europe needs to adapt and play its part by consolidating its core alliances and by creating new ones. In this long-lasting dynamic and thus uncertain context, to resist adverse international political change and other disturbing parameters, resilience must be deeply embedded into a long-term European exploration strategy. Indeed, the inherent feature of space exploration is the longlead time for its preparation, decision-making, development, and operations, that usually lasts over several decades. It is therefore crucial to emphasise that decisions made, or not, at the Council meeting at Ministerial level end of 2022 (CM22) can have impacts into the 2030s and even 2040s. From the onset, the Terrae Novae strategic roadmap has built-in the notion of more European autonomy, leadership and identity. Recent geopolitical events are now fully reinforcing the unavoidability of this approach. Not having autonomous capabilities is indeed a hard lesson learned: developing major scientific instrumentation or technological demonstration capabilities without mastering the delivery to their destination bears a high programmatic and financial risk. Such freedom of action is not incompatible with international cooperation. Being a reliable partner having its own dissimilar redundancy in selected activities is a strong asset. Autonomy and leadership are the prerogative of major economic and political powers that influence the international setting. It is up to our decision-makers to choose to be part of this endeavour, and to further project Europe’s soft power into the Solar System for the benefit of this and the next generations. This document is meant to enlighten such decisions. The scope of this strategy is to create the framework for an actionable and resilient European exploration approach into the next decade that is commensurate with a reasonable ambition.

In 2014, the ESA Council at Ministerial level adopted the “Resolution on Europe’s space exploration strategy”2. The European Exploration Envelope Programme (E3P, branded Terrae Novae in 2021) was created in 2016 to deliver Europe’s space exploration strategy. The programme brings together all ESA’s exploration activities in a single programme. The mission of the Terrae Novae programme is to lead Europe’s human journey into the Solar system using robots as precursors and scouts, and to return the benefits of exploration back to society The Terrae Novae mission statement reflects the raison d’être of the exploration programme and its intrinsic value, i.e., to explore and expand human presence to exploration destinations where humans will one day live and work, making new discoveries and learning about our past and preparing for the future. The mission statement is completed with Europe’s downto-Earth strategic orientations to produce scientific, economic, inspirational, and global cooperation benefits for society. The ESA Council at Ministerial level in 2019 has positioned Europe at the forefront of international exploration campaigns in the 2020s with new capabilities and assets to be deployed for humans in deep cis-lunar space, and for Mars robotic exploration as depicted in and are opening-up new scientific fields while delivering important socio-economic benefits to Europeans. At the beginning of this decade exploration is at an unprecedented crossroads with many established and emerging countries on a course to step-up their posture in exploration, with geopolitical turmoil putting in question historical partnerships and at the same time private investors gaining record momentum and influence. essential capabilities exist, notably in autonomous transportation. Further, all current and planned European assets are contributing to and operated in the context of a partnership. This contrasts with the autonomous transportation capabilities and platforms of International Partners (operational and/or under development). Without new European exploration developments matching this trend, unbalanced dependencies will increase as the global scene develops in the 2030s. capabilities by the end of this decade 3 EL3 (European Large Logistic Lander) is yet to be approved; ESM - European Service Module; I-HAB - International Habitation Module; ESPRIT - European System Providing Refuelling, Infrastructure and Telecommunications; ERO - Earth Return Orbiter; SRL - Sample Return Lander; OS - Orbiting Sample; NET - not earlier than M a r g o E S M A S 2 o m m e rcia l S S S S rio n cre w a te w a y S S e rs e ve ra n ce in cl.la n d e r S S e n is e ryo l u n a 2 2 i Exploration missions maximise scientific value as well as they stretch our imagination and technological capabilities. Exploration contributes to the competitiveness and growth of European industry by pushing the frontiers of knowledge and enabling applications in other fields of the economy. at ESA in key domains, including economy, science, and global cooperation, while also being inspirational and contributing to global challenges. The practical knowledge, products, services, or applications derived from exploration activities have indeed potential to contribute to global challenges, UN Sustainable Development Goals, and pressing European policy priorities by providing solutions in areas such as environment \& climate (responsible consumption, resource management and carbon footprint), up to healthcare \& wellbeing, including novel medical technologies. 3 Lessons learned from 50 years of human and robotic exploration Over 50 years, ESA has a rich history in human and robotic exploration missions, independently and together with its international partners: from Spacelab to Columbus and the Cupola and now I-HAB and ESPRIT; from the ATVs to the ESMs and now the proposed European Large Logistics Lander (EL3), from Mars and Venus Express to Huygens and Rosetta and now ExoMars and the Mars Sample Return campaign. This vast experience has allowed ESA to build up knowhow and lessons learned from which future missions can benefit. By considering the lessons learnt, ESA can become more resilient as well as self-standing when needed or an almost equal partner in the global exploration scheme during the next move. Summary lessons learned for the core Terrae Novae strategy and activities A decision or lack thereof has long-lasting consequences, for example the rejection of the Hermes programme, the discontinuation of the ATV or the missed decision in 2012 of developing capability for precision robotic lunar landings. Making the right move at the right time and nurturing it renders long-term benefits and can adequately position Europe in the forefront for decades to come. A focus on developing complex scientific payloads and technology demonstrators is a high strategic risk if not accompanied by an independent means of launching and landing; this has been exemplified with the ExoMars Rover and the Russian Luna programme cooperation. Alleviating some dependencies on International Partners, and thus becoming more resilient and a leader in exploration thus to attracting new partners into European projects is a clear thread in the Terrae Novae strategic roadmap. Europe, mainly via its industry, has a commensurate role to play in the next generation of service-based LEO infrastructures. The Lunar pressurised mobility and human Mars mission habitats can be a significant strategic move to continue securing the leadership in pressurised modules. Being able to land significant payloads on the Moon would have long lasting benefits for Europe’s role in lunar exploration. The Earth Return Orbiter can be seen as a precursor of a cycler-like logistics provision to support human Mars exploration, or contribution to Human transit capabilities. Building-up expertise in Mars Entry Decent and Landing (EDL) should be continued following the ExoMars experience until end-to-end mastering of such technologies. In the same vein, a second generation of heavier payload EDL technologies should be initiated in order to play a critical role in supporting human exploration. There is an absolute need to elaborate a long-term coherent and encompassing strategic roadmap, by anticipating changes and enabling Europe to be at level playing field with other powers as the acceleration of exploration activities is a given for decades ahead. 4.1 A vision for exploration embedded in an ESA-wide narrative The Terrae Novae 2030+ strategy roadmap unifies the exploration goals at large in an ESA-wide narrative. In times of increasing tensions and moving international context, strategic resilience and European autonomy is paramount and central to this encompassing narrative. The vision of the Terrae Novae 2030+ strategy roadmap is for Europe to step up further and enter the top tier in selected exploration areas with a sustained presence in and utilisation of Low earth Orbit (LEO), the first European astronaut to the Moon surface before 2030 and Europeans to Mars by 2040. Following discussions at the recent Space Summit, it includes considering an option of cargo and crew transportation, entailing the adequate launcher capabilities in an independent manner, for LEO up to the Moon and back. This vision provides Europe in the 2030s the required breadth of advanced capabilities and programme magnitude to play a key role on the global space exploration scene, commensurate with its political and economic weight in the world, still significantly lower than the main actors, namely US and China. In this way, Europe will remain at the forefront of exploration and be part of new and exciting scientific discoveries and technological advancements, secure for its citizens the socio-economic benefits stemming from exploration and play its role as source of inspiration and innovation for a circular economy and contribute to the sustainable development goals. “We want Europe to benefit from space as much as the US and China. We already have the required expertise, knowhow, and industrial capacity. What we need now is a common European space vision and ambition.” ESA Agenda 2025 In line with the ESA Agenda 2025, Terrae Novae 2030+ also aims to stimulate a more vibrant and dynamic commercial space ecosystem. Current service-based space commercialisation initiatives built on governmental investments are shifting towards private investments in the whole value chain, creating a true “NewSpace” shift. Europe cannot afford to miss this transformation and Terrae Novae 2030+ fully embeds and accompanies this objective. In this context, wherever appropriate, Terrae Novae 2030+ encourages a shift from hardwarefocused procurement to implementing activities through purchase of end-to-end services. Leveraging commercial activities of new and established actors in the value chain, including start-ups, SMEs, mid-caps as well as large system integrators. This approach will decrease public sector investments in space infrastructure and open-up development opportunities to the private sector. The Terrae Novae 2030+ vision is translated into the following ambitious goals for each of the exploration destinations (Figure 5): Ensure continuity in LEO by ensuring a continued presence on the ISS until its decommissioning and preparing for post-ISS service-based commercial LEO infrastructures as a primary destination for scientific research and deep space exploration preparation Realise the ambition to have the first European astronaut on the Moon surface before 2030 by providing autonomous Moon landing capabilities for European-led missions including within an international cooperation context, developing scientific and infrastructure assets, and preparing in turn for sustained lunar exploration in the 2030s, possibly also seizing new cooperation opportunities in human landing and surface mobility capabilities Implement a vision for long-term robotic exploration of Mars, that will pave the way for the horizon goal to have the first European to Mars by the end of the next decade by taking leadership in e.g. survivability technologies, mastering radioisotope power sources, entry, descent, and landing for small and eventually large logistics payloads, and to expand scientific knowledge of a sister world The common thread for all three destinations is that autonomous logistics capabilities will allow ESA to take up strategic roles in terms of: European autonomy in end-to-end capabilities – from launch to landing – in order to define and implement Europe’s own science and technology roadmaps resilience of capabilities, providing options to choose between fully European-led activities or interdependent cooperation projects with international partners depending on the context, available resources and benefits evolved partnerships as a means e.g. to offset Europe’s needs for post-ISS LEO utilisation, and to participate in international Lunar surface activities and Mars human transit and surface missions opening options for leadership in terms of capacity building and inviting other partners into specific European activities. evolution of Europe’s role at each of the three destinations in the period 2020-2030 and 20302040. It includes the inevitable commercialisation era of LEO, a cargo and possibly crew capabilities, the implementation of European-led activities at the Moon, and the longer-term opportunities for Mars exploration. It is recognised that the most ambitious goals such as a sustained lunar presence or human Mars exploration will depend on international cooperation. This is one more reason to position Europe well ahead of these fascinating challenges. Terrae Novae in a broader view offers a programmatically diverse, scientifically rich, technologically innovative, and highly inspirational exploration programme that is worthy of our growing European aspirations. By pushing the boundaries, the Terrae Novae 2030+ notional mission roadmap will provide Europe: Directly or indirectly, constant access to LEO for European astronauts, including science opportunities and habitation services and test beds Market shares in an expanding LEO and cis-lunar space economy More autonomy in cargo transportation to and human exploration on the Lunar surface; Increasing scientific knowledge of the Moon and Mars and also of our understanding of Earth Thematic leadership in technology areas e.g. power, surface mobility, habitation or ISRU, to be used across destinations Significant elements on the critical path leading to the horizon goal of Europeans to Mars such as human transit habitation or large robotic Mars surface precision landing Expansion of European excellence in communication, navigation and weather services, with applications from Earth orbit to Moon and Mars Elements of robustness, including a variety of opportunities for small and midsize actors to have visible and meaningful roles on the critical path of the roadmap Opportunities for highly visible European-led missions, enabling new partnerships The know-how and technical competence of industry, combined with the knowledge gained from the European scientific community, form the backbone of the Terrae Novae 2030+ strategy roadmap. Through development of exploration capabilities and implementation of significant exploration missions delivering research and commercial benefits, Europe will be more competent, more competitive and more agile. At the same time, this will reinforce its credibility as a trusted partner in major international projects. The Terrae Novae 2030+ notional high-level mission roadmap for the period 2028-2040 is depicted in Figure 7. It is a bottom-up exercise not constrained by pre-defined top-down budget limitations or purely scientific drivers. However, for affordability reasons, it is not assumed that Europe will develop autonomous European capabilities for human lunar or Mars surface landing systems, or the corollary super-heavy launchers.

The integrated Terrae Novae 2030+ long-term strategy aims to develop sustainable activities in each of the three exploration destinations and to maximise synergies between activities related to LEO, Moon and Mars (cf. Figure 8). European activities and capabilities are prioritised to support as much as possible destinationspecific objectives. As a goal, industrial activities in a specific destination are likely to find synergies in the other destinations. This makes Terrae Novae truly an integrated programme across LEO, Moon and Mars. Examples of synergies for LEO, Moon and Mars objectives Mobility solutions for search for life in the Martian subsurface or prospecting of resources on the Moon: the extensive experience in mobility on Mars (Rosalind Franklin and also Sample Fetch Rover preliminary studies) will significantly contribute to the development of mobility solutions on the Moon. In turn, the evolved mobile systems on Moon will contribute to the Mars scenario Smart solutions for future Mars transit, Moon and Mars surface habitation will build on the extensive experience gained in LEO and continuing in deep space at the Gateway Advanced Mars Entry, Descent and Landing (EDL) demonstration from LEO In the area of complex operations, collecting a sample container in Mars orbit is highly relevant for LEO in-orbit rendezvous, servicing and refuelling technologies. 4.4 Science: the ultimate destination While curiosity is the engine of space exploration, knowledge is the ultimate destination. As science and its applications will make space exploration a greater reality, the knowledge acquired will reveal our history, inform our future, and give us a mirror – like the Pale Blue Dot – for an enhanced understanding of ourselves and our environment. Terrae Novae aims to maximise unique opportunities for performing science to the widest possible European scientific communities in three disciplines: Life Sciences (including space biology and space health research), Physical Sciences, and Moon and Mars Sciences. Terrae Novae science is done on ground, in LEO, around and on the Moon, and Mars and Mars orbit. A comprehensive strategy for science is complementing the Terrae Novae 2030+ strategy roadmap aiming at ensuring continuous opportunities on ground platforms and in all destinations to: perform basic research to understand physical and biological phenomena translate space-acquired knowledge and know-how to support space exploration and address problems on Earth optimise safety, health, and performance of humans in deep space search for life in and enhance the habitability of space. Terrae Novae enables a variety of high-quality science and supporting technology, thematically focused on spotlights that include activities from the three disciplines and abundant opportunities for multidisciplinary and interdisciplinary science activities. Terrae Novae Science Spotlights Humans living on other worlds The support and facilitation of sustainable life on other worlds with focus on Moon and Mars Astronauts 2.0 The review and redefinition of the capabilities, needs, and risks of the next generation of astronauts, recognising that they will represent a broad demographic and will support a range of mission profiles Space travel and transport Science contributing to improve space travel and transport Origin, evolution and protection of extra-terrestrial life Using the Terrae Novae exploration destinations to contribute towards the search for past and present life while safeguarding existing life Exploring the principles of nature using the exploration destinations Making use of the space environment, the Terrae Novae exploration destinations and space analogues to address basic science questions The nature of exploration destinations The formation, evolution, and environmental processes that created and defined the Moon and Mars of today, and what they can tell us about our own planet’s history 5 Terrae Novae 2030+ destination strategies The Terrae Novae 2030+ strategy provides a strong top-down vision and guiding principles. Each destination has its own specific international, economical, technical, and scientific context. Low Earth Orbit Optimise the use of the ISS during its remaining lifetime and prepare human post-ISS activities, including fostering its commercial use and supporting scientific research and the exploration of Moon and Mars The US approach to post-ISS LEO activities is firmly commercial based. Europe will have to adapt to this situation by defining an anchor customer approach whereby opportunities will be given to European industry to provide a service-based offer, ESA being a customer and not an owner of infrastructure. The offer will have to include access to in-orbit infrastructure a well as upload and download capabilities. The European contribution to the “concept of operation” or traffic model of future E activities needs to recognise that the post-ISS infrastructure(s) will have a lifetime of at least 20 years. Providing one or even several pressurised modules, like for Columbus on the ISS, to a commercially driven space station has limited ‘barter’ leverage. ndeed, the “payback” will be short-lived, in the order of only a few years. Conversely, providing up and down mass contributions has the advantage of filling a longlasting need over the whole lifetime of an orbital infrastructure. It will also build on a postATV (Automated Transfer Vehicle) and re-entry know-how to acquire a strategic position for LEO logistics, and in addition initiating a stepping-stone towards a crewed vehicle if, and when the political appetite exists. As in the case of the proposed European Large Logistic Lander (EL3), LEO transportation capabilities for logistics are appealing for European autonomy and resilience. ESA will take this into account as discussions continue with international and foreign industrial partners alongside ESA’s own concept and service studies. The SciHab concept: meeting the European needs for LEO utilisation post-ISS The “Sci ab” Science and abitation concept is proposed as Europe’s future central needs for LEO utilisation (cf. Figure 9). The concept includes a variety of scenarios and levels of ambition. The concept should therefore be regarded as a range of possible needs to Europe’s future challenges and ambitions in LEO. It should NOT be thought of as an institutionally owned classical, stand-alone development project. It is rather a statement of intent for industry to take into account when assessing their technical solutions. It is based on defined European user-needs, relevant and useful in the evolving LEO context (both institutional and commercial). The capability is to be provided and operated on a commercial basis by the private sector. The common denominator of all potential variations of the SciHab concept is to provide access to one (or more) LEO science \& habitation platform(s), with the option to add modular extension elements that might be of interest to Participating States but also commercial users. Combining the SciHab concept with independent transportation capabilities would provide Europe with an opportunity to take up a high level of strategic autonomy, resilience, and leadership in an otherwise fragmented LEO landscape. However, such a capability implies a financial commitment beyond the currently foreseen budget corridor of E3P. The institutional needs of ESA only (1-2 astronauts per year and a few tonnes of cargo) would potentially not warrant the multi-billion Euro investment. Such an investment could more easily be justified if other – non-ESA – users emerge, especially in the commercial domain. The three relevant markets are: commercial research, in-orbit manufacturing of high value goods such as semiconductors or optical fibres, and space tourism and other luxury industry. US commercial space station providers see all three as promising, but they see space tourism as the key market which is already a reality. In addition, new emerging government actors could be interested by accessing LEO for prestige and science reasons, and thus Europe could offer a third way, alongside the US and China. If supported, such a political narrative will significantly foster Europe’s soft power.

Moon Increase European strategic autonomy making ESA a reliable and visible partner in sustainable human and robotic exploration, with a contribution that will bring the first European astronaut on the Moon surface before the end of the decade, secure European scientific discoveries, and prepare the grounds for the next capabilities, technologies and partnerships. build-up of capabilities in cis-lunar space and on the Moon surface. It reflects the European ambition for autonomous roles (in green) but also illustrates how these roles and corresponding capabilities will enable selective contributions necessary to achieve the ultimate goals in a larger partnership (in white). Moon in a single image. It is expected that these capabilities will secure the ESA ambition of “European oots on the Moon”, i.e. European science and technology together with a first European astronaut on the Moon surface as soon as 2030. The Argonaut EL3 lander in the left foreground has delivered a combination of cargo items, scientific payloads, and small robotic assets (rovers). An Artemis crew including a European astronaut are unloading the lander assisted by a robotic arm and preparing lunar surface exploration activities. An earlier Argonaut lander can be seen in the background, supporting several payloads and instruments. In orbit the Gateway including European elements provides the staging post to the lunar surface, as well as other deep space destinations The Moon exploration strategy will prepare Europe to implement strategic autonomy in its lunar exploration activities. A significant science \& utilisation component is an essential and integral part of the strategy. Ongoing Gateway utilisation should be complemented with regular major surface developments, initially synchronised with EL3 flights. Crucially, the research and eventual implementation of ESA’s Space Resources strategy will be embedded in Terrae Novae 2030+ and find its first practical application. Ultimately, this Moon exploration scenario will prepare Europe to take-up a significant role in lunar surface pressurised mobility and/or habitation – perhaps including a permanent research infrastructure. These capabilities will also be essential building blocks for human Mars exploration. Beyond the scientific discoveries the benefits are expected to cover acceleration of technology development, economic footprint of space services, European identity in the geopolitical arena and inspiration for European society and future generations. The near-term decisions in this decade will define the extent to which Europe will be able to benefit in the future by being a significant player.

Mars Undertake robotic precursor missions to continue the search for life, close strategic knowledge gaps and develop capabilities to prepare for human Mars exploration. Following the approval of the Mars Sample Return campaign, Europe has reached one of its long-standing exploration goals, that of embarking on a robotic return of samples from Mars. Europe’s pathway to Mars must now reflect on the next steps, while keeping focused on the horizon goal of the Global Exploration Roadmap, the eventual human exploration of Mars (Figure 12). To meet this challenge, the Terrae Novae Mars scenario is based on a set of principles aiming at building a robust sequence of mainly European-led robotic missions that are open and beneficial to all players in Mars exploration. The notional sequence of flagship Mars missions in Figure 13 comprise a communications and navigation network mission, a weather orbiter and lander network mission, an ice-drilling, ISRU and regenerative fuel cell demo landed mission, and a next-generation precision-landed astrobiology rover mission. The candidate missions are conceived to equip Europe with capabilities in telecommunication, navigation, and climate monitoring (from orbit and the surface) which will enable future ESA and international missions. An emerging new opportunity for post-MSR cooperation is the Mars Life Explorer (MLE) which has been prioritised in the US National Academies of Sciences, Engineering, and Medicine Decadal Survey released in April 2022. Themed candidate missions are complemented by a campaign of small and fast-track ESA-led missions, intended as a regular series of missions offering opportunities for complementary science as well as for focused technology flight demonstrations. While robotic missions to Mars orbit and surface are advancing knowledge and capabilities around and on Mars, the ultimate human journey to Mars is in parallel prepared through synergies with scientific research and technological advancement of human exploration in analogue facilities on Earth, in LEO, around and on the Moon. The strategy seeks to utilise every launch opportunity to Mars when affordable. The notional strategic roadmap of Mars missions is intended to reflect the growing heritage and ambitions of Europe in Mars exploration, consolidating and advancing key technological capabilities that will secure Europe’s independence of action at Mars, closing strategic knowledge gaps. The Mars strategy roadmap following the era of Mars Sample Return will, together with planned activities in LEO and at the Moon, place Europe and its partners in a strong position, by the end of the 2030s, to safely embark on the grand adventure of human Mars exploration with the potential for strong European contributions in areas such as advanced life support systems for transit habitation, logistics transportation, pressurised mobility, ISRU or surface power systems.

In essence, the Terrae Novae 2030+ notional mission roadmap is a sequence of candidate missions designed to deliver the Terrae Novae 2030+ goals. The roadmap has been created as a flexible instrument with options to tune the decisions, in view of scientific and technological breakthroughs and considering the evolving political and programmatic landscape, as well as the level of ambition and affordability at the time of the decisions. The roadmap puts exploration into an ESA-wide perspective, from human space transportation, including future launchers, to utilisation of LEO, up to sustainable Moon and Mars exploration. Overall, the thread of being able to launch and deliver payloads to LEO, the Moon and Mars is a strategic long-term objective, to ensure constant science outcomes and technology developments, assuring Europe a seat at the big table of space explorers. The scenarios in the Terrae Novae 2030+ strategy roadmap for LEO, Moon and Mars, and the integration of those scenarios have been derived bottom up, unconstrained by budget assumptions and purely scientific drivers. The scenarios are the basis for dialogue with Participating States to establish the priorities for the Terrae Novae programme. At CM22, specific decisions will be required to ensure long-term European capabilities (e.g. in LEO) and to prepare the next steps in deep space (e.g. for lunar surface exploration and preparing for humans to Mars). The strategic roadmap work has supported these decisions by informing the selection of new phase A/B1 mission studies, including technology maturation, and giving major orientations for the future evolution of the programme. Importantly, the strategy is providing a narrative needed for political decision makers and taxpayers’ appreciation. Exploration is indeed an investment for future prosperity. t generates high quality jobs and immediate economic return. Exploration science and technologies are a driver and accelerator for sustainable development and have the unique potential to transform into innovative solutions who make life on Earth more productive, clean, and sustainable, securing a safe future for our planet and generations to come. The strategic considerations in this roadmap provide a consolidated proposal to all stakeholders in Europe (governments, space agencies, the science community, and industry including the non-space sector) as well as a message towards our valued international partners that Europe has a direction of travel. It provides a relevant reference document for the work of the High-level Advisory Group on Human Space Exploration for Europe, mandated during the European Space Summit on 16 February in Toulouse. It is now for the political decision-makers to define their level of ambition so that ESA and all its stakeholders can translate this strategy roadmap into reality.

\section{Entidades Nomeadas}

\subsection*{DATE}
\begin{itemize}
\item 50 years (2)
\item June 2022 (2)
\item the beginning of this decade (2)
\item 16 February (1)
\item 16 February 2022 (1)
\item 2012 (1)
\item 2014 (1)
\item 2016 (1)
\item 2019 (1)
\item 2020-2030 (1)
\item 2021 (1)
\item 2022 (1)
\item 2025-2030 (1)
\item 2030s (1)
\item 2040 (1)
\end{itemize}

\subsection*{GPE}
\begin{itemize}
\item Europe (50)
\item US (5)
\item China (3)
\item Toulouse (2)
\item France (1)
\item SRL (1)
\end{itemize}

\subsection*{LOC}
\begin{itemize}
\item Mars (49)
\item LEO (10)
\item Earth (5)
\item Solar System (2)
\item the Moon surface (2)
\item Earth Thematic (1)
\item Gateway (1)
\item Moon (1)
\item Moon and Mars The (1)
\item the Moon ( (1)
\end{itemize}

\subsection*{ORG}
\begin{itemize}
\item ESA (27)
\item Terrae Novae (4)
\item EDL (3)
\item Exploration Science Advisory Committee (3)
\item ISRU (3)
\item ISS (3)
\item post-ISS (3)
\item Advisory Group on Human Space Exploration (2)
\item ESA Council (2)
\item European Large Logistic Lander (2)
\item European Space Agency (2)
\item HESAC (2)
\item International Partners (2)
\item Artemis (1)
\item ATV (1)
\end{itemize}

\subsection*{PERSON}
\begin{itemize}
\item Josef Aschbacher (2)
\item Rosalind Franklin (2)
\item ryo (1)
\end{itemize}

\subsection*{PRODUCT}
\begin{itemize}
\item Agenda 20251 (1)
\item Cupola (1)
\item Orion (1)
\item OS - Orbiting Sample (1)
\item Sample Fetch Rover (1)
\end{itemize}


\section{Observações}

Este documento foi gerado automaticamente a partir da fonte selecionada sobre exploração espacial.
As frases do resumo foram escolhidas com base num modelo de linguagem de n-grams (trigramas) com scoring por log-probabilidade média e suavização de Laplace.
As entidades nomeadas foram extraídas com spaCy (modelo \texttt{en\_core\_web\_lg}) e posteriormente filtradas.

\begin{thebibliography}{9}

\bibitem{source}
ESA. \textit{Terrae Novae 2030+ Strategy Roadmap} 2022. Disponível em: \url{https://esamultimedia.esa.int/docs/HRE/TerraenovaeStrategyRoadmap.pdf}.

\end{thebibliography}

\end{document}
